%% kapa: exam=1 round=21 year=2010 period=2 form=A subject=appraisal_law subject_ko=부동산관계법규 seq=4 number=44
%% source: Q-Net 감정평가사 기출문제 (cst00309) artlSeq=5024181 file=extracted/1008/5024181/2010년 제21회 감정평가사 1차시험 기출문제/2교시.hwp
%% conversion: kordoc 4.24.51 -> tools/gen1.py (auto); review=proofread-gpt-sol; flags=none
\begin{question}{44}
국토의 계획 및 이용에 관한 법령상 용도지역 지정 의제에 관한 설명으로 옳지 않은 것은?
\begin{choices}
  \item 「항만법」에 따른 항만구역으로서 도시지역에 연접한 공유수면은 도시지역으로 결정ㆍ고시된 것으로 본다.
  \item 「산업입지 및 개발에 관한 법률」에 따른 국가산업단지로 지정ㆍ고시된 지역은 도시지역으로 결정ㆍ고시된 것으로 본다.
  \item 「택지개발촉진법」에 따른 택지개발예정지구는 도시지역으로 결정ㆍ고시된 것으로 본다.
  \item 「어촌․어항법」에 따른 어항구역으로서 도시지역에 연접한 공유수면은 도시지역으로 결정․고시된 것으로 본다.
  \item 관리지역의 산림 중「산지관리법」에 따라 보전산지로 지정ㆍ고시된 지역은 자연환경보전지역으로 결정ㆍ고시된 것으로 본다.
\end{choices}
\end{question}
\answer{5}
